\section{Results}\label{sec:results}

\subsection{Accept / reject outcomes}

Every statement in the corpus is run through the full pipeline and the verdict recorded.
Of \CorpusTotal\ statements, \CorpusAccepted\ are \accepted\ and \CorpusRejected\ are
\rejected. A \texttt{field} entry marked \textsuperscript{*} lacks manual-transcription
attestation and does not count towards the collection requirement.

\begin{longtable}{@{}l p{0.44\linewidth} c p{0.22\linewidth}@{}}
\toprule
\textbf{Id} & \textbf{Statement} & \textbf{Verdict} & \textbf{Reason} \\
\midrule
\endhead
\bottomrule
\endfoot
\texttt{demo-001} & \small Combi, on go au kwatt en taco, j'ai les dos pour la course. & \accepted & \small -- \\
\texttt{demo-002} & \small Le réseau ne passe pas, je peux même pas send le message au combi. & \accepted & \small -- \\
\texttt{demo-003} & \small Ils ont encore coupé le courant au kwatt, je wanda ! & \accepted & \small -- \\
\texttt{demo-004} & \small Baisse un peu le prix des tomates, je n'ai pas assez de nkap pour buy tout ça. & \accepted & \small -- \\
\texttt{demo-005} & \small Avec cette pluie, je préfère nang à la piol au lieu de waka dehors. & \accepted & \small -- \\
\texttt{demo-006} & \small Le taximan est parti fala l'essence, sa caisse est à sec. & \accepted & \small -- \\
\texttt{demo-007} & \small Combi, viens tchop le poisson ici, on peut discuter le prix. & \rejected & \small syntax no table entry \\
\texttt{demo-008} & \small Le bendskin va me bring au kwatt, je ne veux pas waka jusqu'à la piol. & \rejected & \small syntax no table entry \\
\texttt{demo-009} & \small On m'a djoss qu'il y a des vols au kwatt, garde bien ton phone. & \accepted & \small -- \\
\texttt{demo-010} & \small Combi, arrête de djoss, on go au school, le prof nous wait déjà. & \accepted & \small -- \\
\texttt{demo-011} & \small Le taximan demande encore les dos, il dit que l'essence coûte trop cher. & \accepted & \small -- \\
\texttt{demo-012} & \small Le courant est back à la piol, mais le réseau ne passe toujours pas. & \accepted & \small -- \\
\texttt{demo-013} & \small Ashia ! & \accepted & \small -- \\
\texttt{demo-014} & \small On dit quoi ? & \accepted & \small -- \\
\texttt{demo-015} & \small Le mbom a tum les beignets au kwatt. & \accepted & \small -- \\
\texttt{demo-016} & \small Je n'ai pas le mbourou, je suis nguémé jusqu'à la fin du mois. & \accepted & \small -- \\
\texttt{demo-017} & \small On va grap au ngola en opep, le ben-sikin est trop cher. & \accepted & \small -- \\
\texttt{demo-018} & \small Le djo va djaffer le soya avec le aboki. & \accepted & \small -- \\
\texttt{demo-019} & \small Ma remé dit que le pater va bring le damé à la locale. & \accepted & \small -- \\
\texttt{demo-020} & \small Le mbéré a hol le man au stationnement. & \accepted & \small -- \\
\texttt{demo-021} & \small Je go fronter au front, le bacho est ndjindja. & \accepted & \small -- \\
\texttt{demo-022} & \small Ashia ! Le djoka est bolè, on pousse déjà. & \accepted & \small -- \\
\end{longtable}


\begin{todobox}[These are demo fixtures, not field data]
Only \CorpusAttested\ of the \CorpusTotal\ statements above are attested field records; the
rest are \texttt{demo} fixtures, as explained in Section~\ref{sec:data}. Re-exporting
against the live corpus database once collection is complete repopulates this table
automatically.
\end{todobox}

\subsection{Why the rejections are rejections}

The two rejected statements are not bugs; they are the documented boundary of the
grammar, and they are more informative than the acceptances.

\subsubsection*{Adverb after the object}

The grammar places adverbs \emph{before} the verb group
(\nonterm{Predicate} $\rightarrow$ \term{ADV} \nonterm{Predicate}) or \emph{inside} it
(\nonterm{VerbTail} $\rightarrow$ \term{ADV} \nonterm{VerbTail}). It has no production
admitting an adverb after a completed object noun phrase, so the parser reaches
\nonterm{Comps} with an \term{ADV} lookahead, finds an empty cell, and reports the expected
set \{\term{NOUN}, \term{DET}, \term{PRON}, \term{PREP}, \term{ADJ}\}.

\subsubsection*{Post-verbal clitic pronoun}

Francanglais permits a clitic pronoun attached after the verb. Modelling this requires
distinguishing a clitic pronoun from a full pronominal noun phrase, and both are currently
\term{PRON}. Adding a post-verbal \term{PRON} alternative to \nonterm{Comps} would collide
with \nonterm{Comps} $\rightarrow$ \nonterm{NP} \nonterm{Comps}, since \term{PRON} is
already in FIRST(\nonterm{NP}) --- an LL(1) conflict.

\begin{limitbox}[The honest trade-off]
Both limitations could be removed by splitting \term{PRON} into \term{PRON} and
\term{CLITIC} at the lexical level and adding the corresponding productions. We did not do
this, because it would mean re-annotating pronoun entries on evidence we have not yet
collected. Reporting two principled rejections is more defensible than widening the grammar
until everything is accepted --- a grammar that accepts every string proves nothing.
\end{limitbox}

\subsection{Verification}

The claims in this report are backed by an automated test suite, run against the same code
that produced every table here.

\begin{center}
\begin{tabular}{@{}lll@{}}
\toprule
\textbf{Check} & \textbf{Scope} & \textbf{Result} \\
\midrule
\texttt{pytest} & 11 test modules, 156 tests & all passing \\
\texttt{ruff check} & full source tree & clean \\
\texttt{pyright} & strict mode, full source tree & 0 errors \\
Playwright e2e & 84 scenarios (24 skipped by platform) & all passing \\
Frontend & TypeScript typecheck and production build & clean \\
\bottomrule
\end{tabular}
\end{center}

Tests that matter specifically for the graded claims:

\begin{itemize}[leftmargin=1.4em, itemsep=0.2em]
  \item \texttt{test\_ll1.py} asserts the constructed table has no conflicting cells --- the
        LL(1) claim in Section~\ref{sec:syntax} is enforced, not asserted in prose.
  \item \texttt{test\_lexer.py} covers multiword matching, the refusal to match across a
        separator, accent folding and homograph preservation.
  \item \texttt{test\_grammar.py} checks that the transformation ledger is non-empty and
        that left recursion is actually gone.
  \item \texttt{test\_corpus.py} checks that a \texttt{field} record without manual
        transcription attestation is rejected at import.
  \item \texttt{test\_demo\_fixture.py} pins the accept/reject verdict of every corpus
        statement, so a grammar change that alters a published result cannot pass silently.
\end{itemize}
